\chapter{Aplicație: electrocardiograma}
\label{ch:ecg}

O bătaie de inimă, văzută pe un ECG, e o secvență de unde cu nume proprii: P
marchează depolarizarea atriilor, complexul QRS depolarizarea rapidă a
ventriculilor, vârful ascuțit pe care oricine îl recunoaște dintr-un ECG,
iar T repolarizarea ventriculilor. Diagnosticul clinic se citește din formă,
din durată și din regularitatea intervalelor dintre bătăi, toate trei cer,
mai întâi, să găsești fiecare complex QRS într-un semnal care aproape
niciodată nu e curat.

\section{De ce wavelets, aici}

Un complex QRS e o tranziție bruscă de amplitudine, o \emph{singularitate}
locală, în limbajul procesării de semnal, suprapusă peste zgomot muscular,
interferență electrică și o linie de bază care alunecă lent. Fourier nu
ajută direct: spune ce frecvențe conține semnalul, nu unde anume apare
tranziția. O transformată wavelet, la o scală apropiată de durata tipică a
unui QRS ($\sim$0.08-0.12s), face exact ce trebuie: coeficienții mari
apar acolo unde semnalul se schimbă brusc, restul rămâne mic. E aceeași
proprietate de localizare din Capitolul~\ref{ch:wavelets}, aplicată la un
semnal 1D real.

Martinez, Almeida, Olmos, Rocha și Laguna~\cite{martinez2004} formalizează
ideea într-un algoritm în trei etape: mai întâi detectează complexele QRS,
apoi delimitează fiecare complex, vârfurile undelor individuale, începutul
și sfârșitul lui, și în final localizează undele P și T. Mecanismul de
bază e detecția de singularități prin maxime de modul ale transformatei
wavelet pe mai multe scări: un wavelet construit ca derivata unei funcții de
netezire produce, lângă o tranziție bruscă, o pereche de extreme de semn
opus, iar tranziția însăși cade aproape de trecerea prin zero dintre ele,
teoria generală de detecție a singularităților din
\cite{mallat1989,mallat2008}, aplicată aici la ECG. Pe bazele de date
standard adnotate manual (MIT-BIH Arrhythmia, QT, European ST-T, CSE),
detectorul de QRS obține o sensibilitate de $99.66\%$ și o predictivitate
pozitivă de $99.56\%$, din o mie de bătăi reale, sub patru sunt ratate sau
inventate.

Miza practică e directă: un ceas care anunță fibrilație atrială citește
un semnal de fotopletismografie (PPG) în locul unui ECG de spital. E lumină
reflectată de piele, modulată de pulsul sangvin, din care extrage
neregularitatea intervalelor dintre bătăi, uneori completat cu o singură
derivație de ECG pe electrozi din carcasă~\cite{perez2019}. Fie că sursa e
un ECG de spital pe douăsprezece derivații sau un semnal de încheietură
zgomotos, primul pas practic e același: găsește bătăile.

\section{Demo}

Demo-ul din tur (Figura~\ref{fig:ecg-demo}) construiește un ECG sintetic ca
sumă de gaussiene pentru undele P, Q, R, S, T ale fiecărei bătăi, o
aproximare standard pentru semnale de test, îi adaugă zgomot muscular
reglabil, îl curăță cu o transformată wavelet 1D (Capitolul~\ref{ch:denoising},
cu pragul universal Donoho-Johnstone, potrivit aici pentru că semnalul e
aproape de ipoteza „neted plus zgomot alb”) și detectează vârfurile R ca
maxim pe fereastra fiecărei bătăi. Frecvența cardiacă iese din mediana
intervalelor R-R. În captura din figură, cu zgomotul la $\sigma=0.25$,
toate cele șase bătăi sunt găsite corect, iar raportul semnal-zgomot urcă de
la $-2.9$ la $2.1$~dB.

\begin{figure}[htbp]
\centering
\includegraphics[width=\textwidth]{35-applications-ecg-demo.png}
\caption{ECG sintetic cu zgomot muscular: denoising wavelet și detecția
  celor șase bătăi, SNR de la $-2.9$ la $2.1$~dB.}
\label{fig:ecg-demo}
\end{figure}
